\documentclass{situationist-preprint}

\title{The Extended Mind Defended}
\msid{MND-09-a}
\msfamily{MND-09}
\msdiscipline{Philosophy of Mind}
\mstype{Principal manuscript}
\mskeywords{extended mind; parity principle; coupling; dispositional belief; cognitive bloat; smartphones}

\begin{document}
\maketitle

\begin{abstract}
Does the mind extend beyond the skull? We argue that it does: cognitive processes extend into external resources that are reliably coupled to the agent, including notebooks and phones. The argument combines the functional individuation of cognitive kinds with a principle of locational neutrality, which together yield the parity principle of Clark and Chalmers, and it fixes the conditions of coupling by four requirements: reliable availability, automatic endorsement, easy access, and past endorsement of the contents by the agent. We show that the four conditions are met by Otto's notebook and by a personal phone used as Otto uses his notebook, and that they exclude the cases that have been urged as reductios: the encyclopedia on the shelf, the internet at large, and a shared document that the agent consults but did not write or adopt. We reply to the coupling-constitution objection, the objection from non-derived content, and the objection that external resources differ in their functional profile from biological memory. The conclusion is that some of our beliefs are partly constituted by marks outside the head, and that some of our cognitive processes run partly through devices.
\end{abstract}

\section{The problem}

\citet{ClarkChalmers1998} introduced the extended mind thesis with two kinds of case. In the first, a subject playing Tetris rotates a shape on the screen to see whether it fits, rather than rotating its image in her head; the rotation on the screen, \citet{KirshMaglio1994} found, is faster than mental rotation and does the same job. In the second, Otto, who has a mild memory impairment, keeps addresses in a notebook he always carries. When he wants to go to the museum, he looks up its address and goes. Inga, whose memory is unimpaired, recalls the address and goes. Clark and Chalmers argued that Otto, like Inga, believed the museum was on 53rd Street before he consulted the notebook, and that his belief is partly constituted by the entry.

The thesis was developed by \citet{Clark2008} and collected with its critics in \citet{Menary2010}. Three objections have shaped the debate. \citet{AdamsAizawa2001,AdamsAizawa2008} argued that the thesis rests on a fallacy, moving from the fact that an external object is causally coupled to a cognitive process to the conclusion that it is part of it, and that cognition requires non-derived content, which marks in a notebook lack. \citet{Rupert2004} argued that external resources differ so much from biological memory in their functional profile, for instance in the absence of interference effects and of generation effects, that the parity between them fails, and that the thesis threatens a bloat of the mind into everything one might consult. \citet{Sprevak2009} argued that functionalism, if it implies the extended mind thesis, implies too much, including that the contents of a library one could consult are among one's beliefs. Others have proposed that the cases are better described as cognition embedded in or scaffolded by an environment, without constitutive extension \citep{Sterelny2010}.

The phone has since replaced the notebook as the standard case. \citet{Chalmers2008} noted that his own phone had taken over parts of his memory, and experimental work has found that people who expect information to remain available externally remember less of the information and more of where to find it \citep{SparrowLiuWegner2011}. Whether such reliance is extension of the mind or dependence on a tool is the question we address.

We defend the thesis, in a form that keeps the endorsement condition of Clark and Chalmers. The condition answers the bloat objection, and it is what distinguishes Otto's notebook from resources that are merely consulted.

\section{Framework}

\begin{definition}[Cognitive kind]\label{d:kind}
A \emph{cognitive kind} is a kind of mental state or process, such as dispositional belief, memory retrieval or mental rotation.
\end{definition}

\begin{principle}[Functional individuation]\label{p:func}
Cognitive kinds are individuated by their functional roles: a state or process belongs to a cognitive kind in virtue of the causal relations in which it stands to inputs, outputs and other states.
\end{principle}

\begin{principle}[Locational neutrality]\label{p:loc}
Whether a state or process plays a given functional role does not depend on where its realizers are located, nor on whether they are biological.
\end{principle}

Principle~\ref{p:func} is the functionalism of MND-08, applied to cognitive kinds in general. Principle~\ref{p:loc} is a consequence of it: a functional role is specified by causal relations, and location and material enter only through their effect on those relations. Together they yield the parity principle of \citet{ClarkChalmers1998}.

\begin{lemma}[Parity]\label{l:parity}
If a process involving an external resource plays the functional role that, were it realized in the head, would make it a process of a cognitive kind, then it is a process of that kind.
\end{lemma}

\begin{proof}
By Principle~\ref{p:func}, membership in the kind is fixed by functional role; by Principle~\ref{p:loc}, the location of the realizer does not affect the role. So a process with the role belongs to the kind wherever its realizers lie.
\end{proof}

The question is then which external resources play the role of internal memory. Clark and Chalmers proposed four conditions, which we adopt as the definition of the relevant coupling.

\begin{definition}[Reliable coupling]\label{d:coupling}
An external resource $R$ is \emph{reliably coupled} to an agent $A$, as a store of information of a given kind, if
\begin{enumerate}[label=(\alph*)]
\item \emph{availability}: $R$ is reliably available to $A$ and typically invoked when the information is relevant;
\item \emph{automatic endorsement}: information retrieved from $R$ is endorsed by $A$ without further scrutiny, as information retrieved from biological memory is;
\item \emph{access}: the information in $R$ is easily accessible to $A$ when required;
\item \emph{past endorsement}: the information in $R$ is there because $A$ put it there or adopted it, having consciously endorsed it at some time.
\end{enumerate}
\end{definition}

Condition (d) is the one Clark and Chalmers marked as the most controversial, and some later defenders of the thesis have dropped it. We keep it, for a reason that comes from the functional role of memory itself. Biological memory stores what the agent has taken in and accepted: its contents are there because of the agent's past cognitive activity, and their being retrieved without scrutiny is warranted by that history. A store whose contents were placed there by others, and which the agent consults but never accepted, does not play that role, however reliable it is. It plays the role of testimony. Condition (d) can be satisfied in more than one way: by writing, by deliberately saving, or by adopting a store whose contents one has reviewed and accepted; what it requires is that the contents entered the store through the agent's own endorsement.

\begin{assumption}[Internal memory]\label{a:mem}
The functional role of dispositional belief stored in biological memory is to provide, on demand and without further scrutiny, information that the agent acquired and accepted, so as to guide reasoning and action.
\end{assumption}

\section{The argument}

\begin{lemma}\label{l:role}
If an external resource is reliably coupled to an agent, its entries, together with the agent's dispositions to consult and act on them, play the role that Assumption~\ref{a:mem} assigns to dispositional belief stored in memory.
\end{lemma}

\begin{proof}
By (a) and (c), the information is available on demand; by (b), it is used without further scrutiny; by (d), it is information the agent acquired and accepted. These are the components of the role in Assumption~\ref{a:mem}. That the information guides reasoning and action is shown by the agent's use of it, as Otto's walk to 53rd Street shows.
\end{proof}

\begin{theorem}\label{t:main}
Dispositional beliefs, and the processes of retrieval that access them, extend into external resources reliably coupled to the agent.
\end{theorem}

\begin{proof}
By Lemma~\ref{l:role}, a reliably coupled resource and the agent's dispositions toward it play the role of dispositional belief stored in memory, and consulting the resource plays the role of retrieval. By Lemma~\ref{l:parity}, they belong to those cognitive kinds. Since the resource lies outside the head, the agent's beliefs and processes of retrieval extend beyond it.
\end{proof}

The same argument applies to processes other than retrieval. When a subject rotates a Tetris shape on the screen to test its fit, or carries a long multiplication on paper, the external manipulation plays the role of an internal transformation of a representation, and the coupling conditions apply with the obvious changes: the device is available, its output is used without further checking, and the representations it transforms are the agent's own.

\begin{corollary}\label{c:cases}
Otto's notebook is part of Otto's memory, and a personal phone used as Otto uses his notebook is part of its owner's memory.
\end{corollary}

\begin{proof}
Otto carries the notebook everywhere, consults it whenever he needs an address, trusts its entries as Inga trusts her memory, and wrote them himself. A phone whose owner carries it, consults its contacts, calendar and notes as a matter of course, trusts what she finds there, and entered or saved those contents herself satisfies (a)--(d) in the same way. Theorem~\ref{t:main} applies to both.
\end{proof}

The corollary is limited to the contents that satisfy the conditions. The notes and contacts on a phone that its owner entered or saved satisfy (d); the encyclopedia and the news feed accessible through the same phone do not. A single device can contain parts of its owner's memory and resources she merely consults.

\section{The bloat objection and the endorsement condition}

\begin{proposition}\label{p:nobloat}
A resource whose contents the agent has not endorsed, however reliably available, is not part of the agent's memory.
\end{proposition}

\begin{proof}
Such a resource fails condition (d). Its contents are not information the agent acquired and accepted, so it does not play the role of Assumption~\ref{a:mem}, and Lemma~\ref{l:parity} gives no ground for counting it as memory. The functional role it does play is that of a source of testimony, which is consulted and trusted without being the agent's own.
\end{proof}

This is the answer to \citet{Rupert2004} and \citet{Sprevak2009}. The encyclopedia on the shelf, the internet, a library: each may be reliably available and easily accessed, and some may be automatically trusted, but none is a store of what the agent has accepted. The same holds for a shared document consulted by many processes that did not write it; PID-16-a uses this point to show that a common catalogue consulted by distinct producers is not a memory that makes them one subject. The endorsement condition is thus not an ad hoc restriction: it separates memory from testimony along the line that functional role already draws.

\section{Objections and replies}

\begin{objection}[Coupling and constitution]
That the notebook is causally coupled to Otto's cognition does not make it part of that cognition \citep{AdamsAizawa2001}. A thermostat is coupled to the room's temperature without being part of it.
\end{objection}

\begin{reply}
The argument does not infer constitution from coupling. It infers membership in a cognitive kind from the possession of that kind's functional role (Lemma~\ref{l:parity}), and it uses coupling only to establish that the role is played (Lemma~\ref{l:role}). The thermostat does not play the role of temperature; Otto's notebook does play the role of his memory. Where the coupling conditions fail, as in Proposition~\ref{p:nobloat}, the argument gives no verdict of extension.
\end{reply}

\begin{objection}[Non-derived content]
Cognitive states have non-derived content. The marks in the notebook have only the content conferred on them by convention and by Otto's interpretation \citep{AdamsAizawa2008}.
\end{objection}

\begin{reply}
The thesis is not that the marks, taken alone, are beliefs, but that the belief is a state of the system of Otto and notebook. Non-derived content is a property of that state, as it is a property of the corresponding state of Inga, whose individual neural components also lack content when taken alone. MND-12 argues that aboutness is primitive and that its distribution is explained by teleosemantics; on that view what bears content is the state with the function of guiding action with respect to the museum, and the system state has it.
\end{reply}

\begin{objection}[Different functional profile]
Biological memory shows interference, generation effects, reconsolidation on retrieval; notebook memory shows none of these \citep{Rupert2004}. So the notebook does not play the role of memory.
\end{objection}

\begin{reply}
These features belong to the fine-grained profile of human biological memory, not to the role that makes a state a dispositional belief. A subject whose biological memory lacked interference effects would not thereby lack beliefs. Assumption~\ref{a:mem} specifies the role at the grain at which folk psychology and the explanation of action individuate belief, and at that grain the notebook and Inga's memory agree. If the finer profile were required, Principle~\ref{p:loc} would be violated in the other direction, since the profile is partly a matter of the biological medium.
\end{reply}

\begin{objection}[Embedding suffices]
Everything the thesis explains is explained equally well by saying that cognition is embedded in or scaffolded by the environment \citep{Sterelny2010}. Extension adds nothing.
\end{objection}

\begin{reply}
The embedded view must say that Otto, before consulting the notebook, did not believe that the museum was on 53rd Street, and so that his walk is explained by a belief formed on reading the entry. Inga's walk is explained by a standing belief. The two explanations then differ at a point where the agents do not: each has a standing disposition, grounded in a store of information each accepted, to head for 53rd Street when the museum comes up. The extended view describes them alike because they are alike in the respect that explanation of action requires.
\end{reply}

\section{What the result does not establish}

The theorem does not show that consciousness extends: the conditions concern dispositional states and processes, not occurrent experience. It does not show that any particular use of a phone extends its owner's mind; many uses fail condition (b) or (d). Nor does it settle whether a group of agents who share reliably coupled resources forms a single subject; the conditions are relative to an agent, and their satisfaction by a shared resource for each of several agents makes it part of each one's memory, not a memory of the group.

\section{Conclusion}

Cognitive kinds are functional, and functional roles are indifferent to the location of their realizers. An external resource that is available, automatically trusted, easily accessed and filled with what the agent has endorsed plays the role of memory, and so is part of the agent's memory. Notebooks and personal phones, used in that way, meet the conditions; encyclopedias, the internet and documents the agent merely consults do not. The mind extends beyond the skull into reliably coupled resources, and the endorsement condition marks where it stops.

\begin{thebibliography}{99}\small
\bibitem[Adams and Aizawa(2001)]{AdamsAizawa2001} Adams, F. and Aizawa, K. (2001). The bounds of cognition. \emph{Philosophical Psychology}, 14(1), 43--64.
\bibitem[Adams and Aizawa(2008)]{AdamsAizawa2008} Adams, F. and Aizawa, K. (2008). \emph{The Bounds of Cognition}. Blackwell.
\bibitem[Chalmers(2008)]{Chalmers2008} Chalmers, D.~J. (2008). Foreword. In A.~Clark, \emph{Supersizing the Mind: Embodiment, Action, and Cognitive Extension}. Oxford University Press.
\bibitem[Clark(2008)]{Clark2008} Clark, A. (2008). \emph{Supersizing the Mind: Embodiment, Action, and Cognitive Extension}. Oxford University Press.
\bibitem[Clark and Chalmers(1998)]{ClarkChalmers1998} Clark, A. and Chalmers, D.~J. (1998). The extended mind. \emph{Analysis}, 58(1), 7--19.
\bibitem[Kirsh and Maglio(1994)]{KirshMaglio1994} Kirsh, D. and Maglio, P. (1994). On distinguishing epistemic from pragmatic action. \emph{Cognitive Science}, 18(4), 513--549.
\bibitem[Menary(2010)]{Menary2010} Menary, R. (ed.) (2010). \emph{The Extended Mind}. MIT Press.
\bibitem[Rupert(2004)]{Rupert2004} Rupert, R.~D. (2004). Challenges to the hypothesis of extended cognition. \emph{The Journal of Philosophy}, 101(8), 389--428.
\bibitem[Sparrow et~al.(2011)]{SparrowLiuWegner2011} Sparrow, B., Liu, J. and Wegner, D.~M. (2011). Google effects on memory: cognitive consequences of having information at our fingertips. \emph{Science}, 333(6043), 776--778.
\bibitem[Sprevak(2009)]{Sprevak2009} Sprevak, M. (2009). Extended cognition and functionalism. \emph{The Journal of Philosophy}, 106(9), 503--527.
\bibitem[Sterelny(2010)]{Sterelny2010} Sterelny, K. (2010). Minds: extended or scaffolded? \emph{Phenomenology and the Cognitive Sciences}, 9(4), 465--481.
\end{thebibliography}

\end{document}
